\documentclass[journal]{IEEEtran}

% Packages
\usepackage[T1]{fontenc}
\usepackage{amsmath,amssymb,amsfonts}
\usepackage{booktabs}
\usepackage{graphicx}
\usepackage{multirow}
\usepackage{array}
\usepackage{url}
\usepackage{balance}

% Custom commands (matching Part I)
\newcommand{\E}{\mathcal{E}}               % Economic decision complex
\newcommand{\Pclass}{\mathcal{P}}          % Projection class
\DeclareMathOperator*{\argmin}{arg\,min}

\begin{document}

\title{Geometric Prediction of Economic Behavior~II:\\Held-Out and Cross-Domain Validation}

\author{Andrew~H.~Bond%
\thanks{A.~H.~Bond is with the Department of Computer Engineering, San Jos\'{e} State University, San Jos\'{e}, CA 95192 USA (e-mail: andrew.bond@sjsu.edu; ORCID: 0009-0003-2599-6158).}%
\thanks{This paper follows ``Geometric Prediction of Economic Behavior: Cross-Domain Validation Across Game Theory and Prospect Theory,'' IEEE Transactions on Computational Social Systems, 2026, doi: 10.1109/TCSS.2026.3734068.}%
\thanks{Manuscript received \today.}}

\markboth{IEEE Transactions on Computational Social Systems}%
{Bond: Geometric Prediction of Economic Behavior~II}

\maketitle

\begin{abstract}
An earlier paper in this journal described economic choices as movements in a space of monetary and non-monetary attributes, priced by one shared distance, and fit it to sixteen aggregate targets.
This paper tests that model on data it was not fit to.
Part I proposed that money should enter in units of the decider's income, and a test registered on the Open Science Framework confirms the proposed coordinate in three countries, while the model that carries Part I's other constants loses to a plain logit.
On 270 risky-choice problems held out from fitting, the model beats expected value and linear random utility and trails cumulative prospect theory by 0.006 in mean absolute error.
A tradeoff shape fit on lotteries predicts play in 110 strategic games, and after a single scale is recalibrated it reaches 97 to 100 percent of a native fit in both directions.
A low-rank form of the shared distance holds on three held-out data sets.
Matched choice problems that standard theories treat as identical move the choices of six language models in the predicted direction for five of seven registered contrasts, while two contrasts fail.
Each passing contrast is also predicted by an established account outside the narrow incumbents.
A fixed temperature fits the observed dose response where Part I's cost-dependent rule failed, in a fit on the same data.
None of these results is from human subjects in a designed experiment, and the human study they size has not been run.
\end{abstract}

\begin{IEEEkeywords}
Computational social systems, decision theory, behavioral prediction, cross-domain transfer, pre-registration, held-out validation, rational inattention, language-model panels.
\end{IEEEkeywords}

%==============================================================================
\section{Introduction}
\label{sec:intro}
%==============================================================================

\IEEEPARstart{A}{n} earlier paper in this journal proposed that economic choices can be read as movements in a space whose coordinates are monetary and non-monetary attributes of an option, such as its consequences for others, its bearing on the decider's identity, and how well its outcomes are known~\cite{bond2026partI}.
One distance, shared across settings, prices every movement, and choice favors the cheaper option.
Fit on bargaining and public-goods targets, the same parameters reproduced several prospect-theory targets.
The model's prediction that ultimatum offers do not depend on the stake failed on high-stakes data from India, and Part I proposed measuring money in units of the decider's income instead.

Part I's evidence was in-sample and aggregate.
Its reviewers asked for individual choices, prediction on held-out data, comparison with competing theories under one protocol, and a closer account of the temperature that turns costs into choice probabilities.
This paper reports those tests.
We call the model Geometric Decision Theory (GDT), and ``the geometric model'' means the same thing.

The first result tests Part I's own proposal.
A registration filed before analysis asked whether money enters the responder's acceptance decision in units of local daily wages, across three countries and stakes up to two hundred days of wages.
The coordinate passes every test registered for it.
The geometric responder built on it, carrying Part I's other constants, is beaten by a plain logit on every held-out country.

The second result concerns prediction on unseen problems.
On 270 risky-choice problems from the CPC18 competition, the model predicts held-out problems better than expected value and linear random utility, and worse than cumulative prospect theory by a small and consistent margin.
A tradeoff shape fit only on lotteries predicts play in 110 published two-player games.
Transferring only that shape, with one scale recalibrated, reaches 97 to 100 percent of a native fit in both directions.
The same construction transfers a self-versus-other tradeoff between budget-line giving and dictator games.

The third result concerns the structure of the shared distance.
Part I assumed it had no interactions between coordinates.
A low-rank form with interactions beats that assumption on three held-out data sets.

The fourth result is a designed test.
Two choice problems with identical payoffs and probabilities, differing only in a non-monetary description, must receive identical predictions from any theory that reads only payoffs and stated probabilities.
GDT predicts the direction in which choices move.
Run on six language models conditioned on 36 personas, five of seven registered contrasts move in the predicted direction and two do not.
Each of the five is also predicted by an established account, for example ambiguity aversion or concern for third parties, so these contrasts reject the narrow incumbents without separating GDT from those accounts.

The fifth result concerns the temperature.
Part I's temperature grew with the cost difference.
On the panel's dose-response contrast it predicted an inverted U where the data rise monotonically, and a fixed temperature, which rational inattention reads as a price of information, fits better on the same data.

Section~\ref{sec:recap} restates the model and how the predictions were frozen.
Sections~\ref{sec:income} to~\ref{sec:temperature} report the five results.
Section~\ref{sec:corroboration} reports related findings in the social-preference literature, Section~\ref{sec:rejected} reports conjectures that failed, and Section~\ref{sec:discussion} states what the results do and do not establish.

%==============================================================================
\section{The Model, the Competing Class, and How Predictions Were Frozen}
\label{sec:recap}
%==============================================================================

We reuse Part I's model~\cite{bond2026partI}.
An option is a displacement $\Delta\mathbf{a}\in\mathbb{R}^9$ from a reference point, with coordinates for monetary consequences ($d_1$), rights ($d_2$), fairness ($d_3$), autonomy ($d_4$), privacy and trust ($d_5$), social impact ($d_6$), virtue and identity ($d_7$), legitimacy ($d_8$), and epistemic status ($d_9$).
Its cost is the Mahalanobis distance
\begin{equation}
    C(\Delta\mathbf{a}) = \sqrt{\Delta\mathbf{a}^{\top}\, \Sigma^{-1}\, \Delta\mathbf{a}},
    \label{eq:mahalanobis}
\end{equation}
and an option $A$ is chosen from a menu with probability
\begin{equation}
    P(A) = \frac{\exp(-C_A/T)}{\sum_{B}\exp(-C_B/T)}
    \label{eq:softmax}
\end{equation}
at temperature $T>0$.
Part I used $T(\delta)=\max(0.5,\,0.24\,\delta^{2.13})$ with $\delta=|C_A-C_B|$, and a diagonal $\Sigma$ with active coordinates $d_6$, $d_7$ and $d_9$.
Sections~\ref{sec:lowrank} and~\ref{sec:temperature} revisit both choices.

The competing class $\Pclass$ contains every model whose choice probabilities depend on an option only through a statistic $\pi(x)$ that omits some coordinates.
Expected utility reads money, payoff-based game theory reads own payoff, inequality aversion reads own and others' payoffs, and cumulative prospect theory (CPT) reads outcomes and stated probabilities.
GDT also reduces each option to one number before choosing, so the contrast is not between scalar and vector models.
It is between models that drop a coordinate before choosing and a model that keeps it.
A model in $\Pclass$ whose statistic omits a coordinate $z$ predicts the same choices for any two options that differ only in $z$.
Section~\ref{sec:projgap} uses this.

\textbf{How the predictions were frozen.}
The income-units test of Section~\ref{sec:income} was registered on the Open Science Framework (\texttt{osf.io/ead4x}) on 2026-09-10, with a signed repository tag on the frozen files, before the analysis ran on 2026-09-12.
The language-model panel (\texttt{prereg-v1}), its revision (\texttt{prereg-v2}), and the low-rank claims (\texttt{prereg-sigma-v1}) were written down and committed to a local repository on 2026-07-04 before their data were analyzed.
The panel's predictions were committed at 09:37 and its results at 14:27 Pacific time that day.
The repository became public on 2026-07-14, after those runs.
The public copies of the frozen files are byte-identical to the committed ones, so a reader can check what was predicted, but the order of prediction and data rests on the author's commit timestamps.
The \texttt{prereg-sigma-v1} claims have no hash file.
We use ``registered'' for the September study and ``frozen before analysis'' for the July ones.

%==============================================================================
\section{Money Enters Through Income Units, as Part I Proposed}
\label{sec:income}
%==============================================================================

Part I encoded money as a share of the stake, so its calibrated model predicted that ultimatum behavior does not depend on the stake.
Offers in the Andersen et al.\ study in Meghalaya, India, fell from 24 to 12 percent of the pie as the stake rose to about a year's income~\cite{andersen2011}, and Part I reported the prediction as failed.
Part I then proposed a monetary coordinate scaled by the stake relative to local income, and stated that it was untested.
This section reports the test.

\textbf{Data.}
The replication package of~\cite{andersen2011} contains 1{,}456 responder decisions from three studies, 820 from Slonim and Roth in Slovakia~\cite{slonim1998}, 178 from Cameron in Indonesia~\cite{cameron1999}, and 458 from Andersen et al.\ in India.
Stakes range from a quarter of a day's wage to 200 days of wages.

\textbf{What was known before registering.}
At comparable stakes in days of wages, mean offers differ across the three populations by 20 percentage points or more.
The strongest form of Part I's proposal, that offers depend on the stake only through its ratio to income, was therefore false before any model was fit, and the registration records it as falsified in advance.
The registered predictions concern the responder.

\textbf{Registered predictions and results.}
Each model is a logistic regression for acceptance with a study-specific intercept and offer-share slope.
P1 asks whether the log of the offered amount in days of wages, with one coefficient shared by all three studies, improves the fit.
It does, with coefficient $0.817$ (standard error $0.147$) and a likelihood-ratio $p=2.3\times10^{-10}$ against a threshold of $0.01$.
P2 asks whether separate coefficients per study fit better.
They do not ($p=0.19$), which is consistent with one shared coordinate but does not demonstrate it, because the registered simulation checked the test's size and not its power.
The ratio of the largest to the smallest per-study coefficient is 2.32, with a bootstrap interval from 1.13 to 9.70.
P3 asks whether the nominal stake matters once the offered amount is known.
Its coefficient is 0.165 per tenfold increase ($p=0.46$), inside the registered bound of 0.25.
P4 fits the shared model on two studies and predicts the third.
On every held-out study its log-loss is within 0.003 nats per decision of a model fit to that study itself, against a registered margin of 0.02.
P6 asks whether proposers in Slovakia and Indonesia change their offers with the stake, and their slopes stay within the registered 1.5 percentage points per tenfold increase.

P5 asked whether the geometric responder, with Part I's encoding and every constant fixed except the monetary variance and the temperature, predicts as well as the logit.
It does not.
Its held-out log-loss exceeds the logit's by 0.033 nats on Slovakia, 0.060 on Indonesia and 0.325 on India, against a margin of 0.03.
The monetary variance went to the upper bound the registration allowed, which reduces the model to Part I's responder, so the income rescaling did no work inside it.

The coordinate Part I proposed passes its tests, and the geometric model built from Part I's other constants fails, most of all on the highest stakes.
The study tested a monetary coordinate and one specification of the model, and it locates the stake effect on the responder side in two of the three populations.

%==============================================================================
\section{Prediction on Held-Out Problems and Transfer Across Domains}
\label{sec:transfer}
%==============================================================================

\subsection{Shared encoding}
To place lotteries and games in one space we encode each option by its expected value and dispersion.
A lottery option $(H,p,L)$ has $\mathrm{EV}=pH+(1-p)L$ and $\mathrm{SD}=\sqrt{p(H-\mathrm{EV})^2+(1-p)(L-\mathrm{EV})^2}$.
A row strategy with payoff matrix $M$, under a uniform belief over the opponent's actions, has $\mathrm{EV}_s=\operatorname{mean}_j M[s,j]$ and $\mathrm{SD}_s=\operatorname{std}_j M[s,j]$, so payoff variance stands in for strategic uncertainty.
Payoffs are scale-normalized per data set.
The cost is the Mahalanobis distance from the ideal of high expected value and low dispersion, and the same $(\sigma_{\mathrm{ev}},\sigma_{\mathrm{sd}},T)$ govern both domains.

\subsection{Data}
Lotteries are the CPC18 description problems~\cite{erev2017,cpc18data}, 270 problems and 27{,}780 first-trial choices made without feedback.
Games are the behavioral game theory data pool of~\cite{wright2017}, 110 two-player normal-form games from eight published studies with human action frequencies.

\subsection{The model beats expected value and random utility on unseen problems and trails CPT}
Each model predicts the choice rate of each problem.
Parameters are fit on a random 70 percent of problems and scored on the other 30 percent, and this is repeated for 20 random splits.
Every model uses a single fitted temperature on each split, which for the geometric model means a fixed temperature rather than Part I's rule.
Table~\ref{tab:cpc18} reports held-out mean absolute error and Table~\ref{tab:cpc18paired} the paired differences.

\begin{table}[!t]
\renewcommand{\arraystretch}{1.15}
\caption{Held-Out Choice Prediction on CPC18 Over 20 Random 70/30 Splits by Problem. Lower Is Better}
\label{tab:cpc18}
\centering
\begin{tabular}{l r r}
\toprule
\textbf{Model} & \textbf{Parameters} & \textbf{Held-out MAE} \\
\midrule
Constant                 & 1 & $0.199 \pm 0.012$ \\
Expected value (logit)   & 1 & $0.149 \pm 0.008$ \\
Linear random utility    & 4 & $0.128 \pm 0.010$ \\
Geometric (GDT)          & 4 & $0.124 \pm 0.009$ \\
CPT                      & 5 & $0.118 \pm 0.008$ \\
\bottomrule
\end{tabular}
\begin{flushleft}
\footnotesize{Mean and standard deviation across splits. Every model is scored on the same splits, so these columns overstate the uncertainty of a difference, and Table~\ref{tab:cpc18paired} gives the paired comparison.}
\end{flushleft}
\end{table}

\begin{table}[!t]
\renewcommand{\arraystretch}{1.15}
\caption{Paired Differences in Held-Out MAE on the Same 20 Splits. Negative Favors GDT}
\label{tab:cpc18paired}
\centering
\begin{tabular}{l r c r}
\toprule
\textbf{Contrast} & \textbf{Mean diff.} & \textbf{95\% CI} & \textbf{Splits won} \\
\midrule
GDT $-$ EV  & $-0.0253$ & $[-0.0284,\,-0.0222]$ & $20/20$ \\
GDT $-$ RUM & $-0.0042$ & $[-0.0058,\,-0.0026]$ & $18/20$ \\
GDT $-$ CPT & $+0.0059$ & $[+0.0035,\,+0.0083]$ & $2/20$ \\
\bottomrule
\end{tabular}
\begin{flushleft}
\footnotesize{Student-$t$ intervals on the 20 paired differences. All three contrasts also reject at $p<10^{-3}$ under a Wilcoxon signed-rank test, and the ordering is the same on held-out negative log-likelihood. The splits overlap, so these intervals are approximate. No multiplicity correction is applied.}
\end{flushleft}
\end{table}

The geometric model beats expected value on every split and linear random utility on 18 of 20.
CPT beats it by about 0.006 in mean absolute error on 18 of 20 splits, with one more parameter.
The CPC18 competition itself was won by models that combine behavioral mechanisms with machine learning~\cite{plonsky2019}, and they predict this benchmark better than CPT.
We compare against decision theories with four or five parameters and make no claim against those forecasting models.

\subsection{A metric fit on lotteries predicts play in games}
We fit $(\sigma_{\mathrm{ev}},\sigma_{\mathrm{sd}},T)$ on one domain and predict the other with no refit, scoring by multinomial negative log-likelihood (NLL) per choice.
The share of the chance-to-native distance closed is $(\mathrm{NLL}_{\text{chance}}-\mathrm{NLL}_{\text{transfer}})/(\mathrm{NLL}_{\text{chance}}-\mathrm{NLL}_{\text{native}})$, so 100 percent means the transfer equals a native fit and a negative value means it is worse than chance.

\begin{table}[!t]
\renewcommand{\arraystretch}{1.15}
\caption{Transfer Across Domains With No Refit. Multinomial NLL Per Choice, Lower Is Better}
\label{tab:transfer}
\centering
\begin{tabular}{l l r}
\toprule
\textbf{Target} & \textbf{Model} & \textbf{NLL} \\
\midrule
Games      & fit on games (native)          & $0.876$ \quad (chance $1.040$) \\
Games      & transfer lotteries$\to$games   & $0.931$ \\
Lotteries  & fit on lotteries (native)      & $0.624$ \quad (chance $0.693$) \\
Lotteries  & transfer games$\to$lotteries   & $0.783$ \\
\bottomrule
\end{tabular}
\begin{flushleft}
\footnotesize{The lottery-fit metric closes 66 percent of the chance-to-native distance on games it never saw. The reverse transfer is worse than chance.}
\end{flushleft}
\end{table}

A metric that has seen only lotteries predicts strategic play two thirds of the way from chance to a native fit (Table~\ref{tab:transfer}).
CPT has no representation of a game and Nash equilibrium none of a lottery, so neither can be scored on this transfer.
The reverse direction is worse than chance, because the smaller game calibration is overconfident on lotteries.

\subsection{What transfers is the shape of the tradeoff}
The cost factors as
\begin{equation}
    C = r\cdot\sqrt{(\cos\theta/\sigma_{\mathrm{ev}})^2 + (\sin\theta/\sigma_{\mathrm{sd}})^2},
    \label{eq:polar}
\end{equation}
a radius $r$ times an angular term fixed by the ratio $\sigma_{\mathrm{ev}}\!:\!\sigma_{\mathrm{sd}}$, which we call the aversion angle.
Choice probabilities depend on $\Sigma$ and $T$ only through $\Sigma^{-1}/T^2$, so the shape of the metric is identified and its overall scale is not separable from the temperature~\cite{bond2026theory}.
Transferring the angle and refitting one scale per domain transfers exactly the identified part.

\begin{table}[!t]
\renewcommand{\arraystretch}{1.15}
\caption{Transferring Only the Angle Recovers Both Directions. Percent of Chance-to-Native Distance Closed}
\label{tab:polar}
\centering
\begin{tabular}{l r r}
\toprule
\textbf{Direction} & \textbf{Full transfer} & \textbf{Angle only} \\
\midrule
lotteries $\to$ games & $0.931$ \; ($66\%$) & $0.876$ \; ($100\%$) \\
games $\to$ lotteries & $0.783$ \; ($<0$)   & $0.627$ \; ($97\%$) \\
\bottomrule
\end{tabular}
\begin{flushleft}
\footnotesize{One scale parameter is refit per domain. The fitted ratios $\sigma_{\mathrm{ev}}/\sigma_{\mathrm{sd}}$ are 0.105 for lotteries and 0.146 for games, both weighted toward expected value.}
\end{flushleft}
\end{table}

With only the angle transferred, both directions reach 97 to 100 percent of a native fit (Table~\ref{tab:polar}).
This analysis was chosen after the asymmetry of Table~\ref{tab:transfer} was seen, so it is exploratory, and its protection is that only one scale is refit.
A structural search over model forms, scored by worst-case transfer in both directions, prefers a variance dispersion feature and a city-block norm to the default.
That search chose its winner on the same two data sets, so it is model selection and not held-out confirmation.

\subsection{The self-versus-other tradeoff transfers between giving tasks}
We restrict the metric to own payoff $s$ and other's payoff $o$, with the ideal $(\max s,\max o)$ in each menu and cost $C=\sqrt{((s^*-s)/\sigma_s)^2+((o^*-o)/\sigma_o)^2}$, regularized by $\sigma\geq0.05$.
Domain A is the budget-line giving experiment of Fisman, Kariv and Markovits~\cite{fisman2007}, 3{,}800 choices by 76 subjects.
Domain B is the seven two-person dictator games of Charness and Rabin~\cite{charness2002}.

\begin{table}[!t]
\renewcommand{\arraystretch}{1.15}
\caption{Self-Versus-Other Transfer. Percent of Chance-to-Native Distance Closed}
\label{tab:social}
\centering
\begin{tabular}{l r r}
\toprule
\textbf{Direction} & \textbf{Full transfer} & \textbf{Angle only} \\
\midrule
FKM giving $\to$ Charness--Rabin & $66\%$ & $100\%$ \\
Charness--Rabin $\to$ FKM giving & $44\%$ & $45\%$ \\
\bottomrule
\end{tabular}
\begin{flushleft}
\footnotesize{Fit separately, the two tasks' tradeoff shapes have cosine 0.96. The larger data set transfers fully and the seven-game fit does not.}
\end{flushleft}
\end{table}

Fit separately, the two tasks' tradeoff shapes have cosine 0.96, and the giving metric predicts dictator play at 100 percent of a native fit with one scale refit (Table~\ref{tab:social}).
The seven-game fit predicts the 3{,}800 giving choices at 45 percent, and transferring only the angle does not improve it, because seven aggregate games do not pin the angle.
Roughly 40 percent of the giving subjects keep everything, which drives one axis to the regularization floor, so the level of other-regard differs between the populations even where the shape agrees.

Two further limits apply to this section.
Iterated best-response beliefs make the game encoding worse, and the uniform belief predicts best, as in~\cite{wright2017}.
The two-outcome approximation for multi-outcome lotteries and the per-data-set scale normalization apply to every model alike.

%==============================================================================
\section{A Low-Rank Metric Beats the Diagonal on Three Held-Out Legs}
\label{sec:lowrank}
%==============================================================================

Part I restricted $\Sigma$ to be diagonal.
A search on lotteries and games suggested a precision matrix of rank one or two would match the full metric and beat the diagonal.
Four claims were frozen before the held-out data were analyzed (\texttt{prereg-sigma-v1}, Section~\ref{sec:recap}).
H1 states that rank one fits within 0.02 NLL of the full metric.
H2 states that the metric transfers at least 85 percent of the chance-to-native distance.
H3 states that rank one beats the diagonal at equal parameter count.
H4 states that rank two adds little over rank one.

\begin{table}[!t]
\renewcommand{\arraystretch}{1.2}
\caption{Low-Rank Claims on Three Held-Out Legs}
\label{tab:lowrank}
\centering
\begin{tabular}{@{}p{0.24\columnwidth} c c c p{0.20\columnwidth}@{}}
\toprule
\textbf{Held-out leg} & \textbf{H1} & \textbf{H3} & \textbf{H4} & \textbf{H2 (transfer)} \\
\midrule
Ultimatum responders, Fraser and Nettle~\cite{fraser2020} & Pass & Pass & Pass & $99\%$ across hunger conditions \\
Prospect-theory items in 19 countries, Ruggeri et al.~\cite{ruggeri2020} & Pass$^{*}$ & Pass & Pass & $99\%$ across countries \\
Six languages & --- & --- & --- & axis cosine 0.818 vs.\ placebo 0.604 \\
\bottomrule
\end{tabular}
\begin{flushleft}
\footnotesize{For H3, rank one against the diagonal gives NLL 0.290 against 0.335 for Fraser and Nettle and 0.667 against 0.680 for Ruggeri et al. $^{*}$Ruggeri et al.\ rank one sits at the 0.02 line, and rank two is the closer low-rank fit. The language leg tests whether coordinate directions survive translation and is described in the text.}
\end{flushleft}
\end{table}

Where the claims can be tested, they hold (Table~\ref{tab:lowrank}).
Rank one or two suffices, the diagonal is beaten, and the metric transfers at about 99 percent across hunger conditions and across 19 countries.
In each leg the precision matrix is fit on five encoded features of the options, not on the nine coordinates.
For the Ruggeri et al.\ items the features are expected value, standard deviation, skew, worst outcome and probability of a favorable outcome, and the Fraser and Nettle leg uses five analogous allocation features.
The fitted loadings were not saved.

\textbf{Language leg.}
The panel's contrast texts were translated from English into Spanish, Chinese, Arabic, Hindi and Swahili by a language model, and a translation was kept only if its LaBSE sentence embedding had cosine at least 0.75 with the English source.
For each contrast, the coordinate axis is the mean embedding of the high-pole texts minus the mean of the low-pole texts in normalized LaBSE space, and invariance is the cosine between a language's axis and the English axis.
Real contrasts give a mean cosine of 0.818 over 35 language-contrast pairs, against 0.604 over 10 placebo pairs.
A behavioral panel ran the 13 contrasts in all six languages, 5{,}616 choices from six models and 12 personas.
The vectors of choice shifts over the seven real contrasts correlate across the 15 language pairs with mean Pearson $r=0.737$, against 0.008 when contrasts are shuffled.

The low-rank result holds for gains.
Under losses the encoding does not capture the reversal of risk attitudes, and transfer from gains to losses on the Ruggeri et al.\ items is zero.

%==============================================================================
\section{Matched Problems Move Language-Model Choices on Five of Seven Contrasts}
\label{sec:projgap}
%==============================================================================

\subsection{Design}
A matched pair $(x,x')$ has identical monetary payoffs and identical lottery structure and differs only in the description of one non-monetary coordinate, here social impact, identity or epistemic status.
Every model in $\Pclass$ whose statistic omits that coordinate predicts identical choices on the two problems.
GDT predicts a signed difference $\Delta=P(A\mid x_{\text{hi}})-P(A\mid x_{\text{lo}})$.
The design, seven real contrasts, two wording placebos with predicted $\Delta=0$, and a four-level dose series were frozen before the panel was run (\texttt{prereg-v1}, Section~\ref{sec:recap}).

Two facts, proved in the companion theory paper~\cite{bond2026theory}, govern what such a design can show.
First, a model's choices depend on an option only through $\pi$ exactly when no pair matched under $\pi$ produces a difference, and a pair matched under several incumbents at once exists exactly when the manipulated coordinate is not a function of their statistics.
A finite set of null results therefore fails to refute $\pi$ without confirming it.
Second, when the manipulation moves one option further from the reference on a coordinate whose cross-terms with the option's other displacements are bounded by its own diagonal term, that option's cost rises and its choice probability falls, for any number of options.
The direction is the same under every strictly increasing recoding of the coordinate, so a confirmed direction is not an artifact of the coding scale.
The size of the effect is not invariant.

\subsection{The panel}
The panel used six models served under the aliases gemma-small, qwen3-small, gemma, gpt-oss, glm-5 and qwen3 by the National Research Platform's model endpoint (\texttt{ellm.nrp-nautilus.io}), at sampling temperature 0.9.
The endpoint resolves these aliases on the server and neither repository records the upstream model versions, so the panel cannot be rerun on identical weights.
Each model answered under 36 personas, crossing risk attitude (averse, neutral, seeking), social orientation (self-interested, fair-minded, altruistic), culture (individualist, collectivist) and numeracy (low, high).
The system prompt read ``You are a human participant in an economics study. Decide as a real person who is \{description\}. Answer honestly and consistently with that disposition.''
Each problem ended with the instruction to answer with only the line \texttt{FINAL CHOICE=<A or B>}.
The panel returned 83{,}682 valid choices from 216 model-persona subjects, 99.3 percent of attempts.

\textbf{Unit of analysis.}
For each subject and contrast, $\Delta$ is the difference between that subject's choice rates on the two problems, and each subject's $\Delta$ is corrected by subtracting the same subject's $\Delta$ on the wording placebo of the same frame.
The test is a Wilcoxon signed-rank test over the 216 subjects (202 for the identity contrast), Cohen's $d$ is the mean over the standard deviation of the subject values, and false-discovery control is Benjamini--Hochberg at $q=0.05$ over the seven real contrasts.
Personas within a model share weights, so the six models are the closest thing to independent units, and the per-model sign count is the primary cross-model statistic.

\begin{table}[!t]
\renewcommand{\arraystretch}{1.15}
\caption{All Seven Registered Contrasts. Five Hold After Placebo Correction and Two Fail}
\label{tab:projgap}
\centering
\footnotesize
\begin{tabular}{@{}l l r r c c@{}}
\toprule
\textbf{Contrast} & \textbf{Coordinate} & \textbf{$\Delta$} & \textbf{$d$} & \textbf{Raw} & \textbf{Corr.} \\
\midrule
A3.1 trust game     & epistemic & $+0.46$ & $0.91$ & 6/6 & 6/6 \\
A1.2 public good    & social    & $+0.32$ & $0.81$ & 6/6 & 6/6 \\
A2.1 audit report   & identity  & $+0.25$ & $0.64$ & 6/6 & 6/6 \\
B1.1 lottery, gain  & epistemic & $+0.17$ & $0.52$ & 6/6 & 6/6 \\
A1.1 gift to shelter & social   & $+0.17$ & $0.56$ & 6/6 & 5/6 \\
\midrule
B2.1 lottery, shared & social   & $-0.01$ & $-0.03$ & 4/6 & 4/6 \\
B1L.1 lottery, loss & epistemic & $+0.06$ & $0.20$ & 1/6 & 3/6 \\
\bottomrule
\end{tabular}
\begin{flushleft}
\footnotesize{$\Delta$ is the placebo-corrected change in choice probability and $d$ is Cohen's $d$. The last two columns count models whose mean change has the predicted sign, before and after placebo correction. A-contrasts are games and B-contrasts are lotteries. The predicted sign is positive for all rows except B1L.1, where it is negative.}
\end{flushleft}
\end{table}

\subsection{Results}
Five of seven contrasts move in the predicted direction after placebo correction and pass false-discovery control (Table~\ref{tab:projgap}).
Four of them have the predicted sign in all six models after correction, and the gift-to-shelter contrast in five, because the smallest model reverses once its placebo is subtracted.
The wording placebos are null.
Four of the five passing contrasts are games.
Of the three lottery contrasts only one passes.
B2.1, which shares a lottery's payoff with the decider's community, shows a raw shift of $+0.021$ that disappears after placebo correction ($\Delta=-0.006$, $p=0.34$).
B1L.1 predicted that vagueness about the odds would reverse its effect under losses, and the observed shift has the gain-domain sign, significant in the wrong direction ($p<10^{-3}$).
The dose series rises monotonically in the pooled data, with its peak at the fourth of four levels where the frozen model predicted the second, and three of the six models are individually monotone.
Magnitudes fall outside the frozen model's predicted range for all seven contrasts, and the trust-game shift of $+0.46$ is about five times the predicted $+0.09$.

\subsection{What these contrasts discriminate}
Each passing contrast holds payoffs and stated probabilities fixed, so each rejects the incumbents as specified over those quantities.
Each also has an established account that predicts its sign.
Vague odds lower the choice of a lottery, which is ambiguity aversion~\cite{ellsberg1961}.
A verified record that a trust-game partner returns half raises sending, which expected utility predicts once beliefs about the partner change.
Benefits to a shelter or a village raise giving, which altruism, warm glow and concern for total surplus predict~\cite{andreoni1990,charness2002}.
An auditor's professional identity raises honest reporting, which identity utility predicts~\cite{akerlof2000}.
These contrasts therefore do not separate GDT from those accounts.
The claim specific to GDT is that one metric with one set of active coordinates predicts the sign of all of them, and the two failed contrasts are failures of that claim.

The panel is a population of language models, not people.
Its results show that the frozen metric predicts the direction of most matched-problem shifts in that population, and they supply effect sizes for a human study.
They are not evidence about human choice.

%==============================================================================
\section{A Fixed Temperature Fits the Dose Response That Part I's Rule Missed}
\label{sec:temperature}
%==============================================================================

Part I's temperature grows faster than the cost difference, with exponent 2.13 fit on aggregate targets.
Fit on CPC18 individual choices, the exponent is 0.42.
An exponent above one makes small cost differences wash out, which predicts an inverted-U dose curve and shrinks effect sizes.

A fixed temperature has a reading from rational inattention.
An agent paying a constant cost per unit of information chooses by a logit whose temperature is that cost~\cite{sims2003,matejka2015}.
The standard multinomial logit follows when the options are exchangeable before information is acquired, and the general solution weights each option by its unconditional choice probability.

We compared three temperature rules on the language-model dry run of the advice-letter instrument built for the human study, 16{,}688 valid choices from 216 subjects in 26 contrast-by-pole cells.
The fit is in-sample on all cells.
The cost-dependent rule has the worst Bayesian information criterion, 124.4 above a fixed temperature and 141.6 above a rational-inattention rule with a fitted prior, which is itself 17 below the fixed temperature.
Across the four dose levels the observed shifts rise monotonically, Spearman correlation $+1.0$, and so do the predictions of both fixed-price rules, while the cost-dependent rule predicts $-0.95$.
These comparisons show which rule fits these data and do not test any rule on held-out data.
The revised rule is part of \texttt{prereg-v2} for the human study.
Two discrepancies remain after the change, the absolute size of the active variances and the encoding of losses.

%==============================================================================
\section{Related Findings in the Social-Preference Literature}
\label{sec:corroboration}
%==============================================================================

\textbf{Low dimensionality.}
In the Fisman--Kariv--Markovits giving data~\cite{fisman2007}, 96 percent of the 76 subjects have convex, downward-sloping giving demand.
A two-parameter constant-elasticity utility fits each subject with mean $R^2=0.54$, and the median weight on own payoff is 0.78.
Finite-mixture and clustering studies of distributional preferences recover about three types, altruistic, inequality-averse and mostly selfish, across large Swiss, German, Danish and U.S.\ samples~\cite{bruhin2019,fehrCharness2025}.
Heterogeneous social behavior thus reduces to a few latent dimensions, as the low-rank result does for risk.

\textbf{One asymmetric axis in social choice.}
Fehr--Schmidt utility is kinked at equality and weights disadvantageous inequality more than advantageous inequality~\cite{fehr1999}.
Our fit to the Charness--Rabin dictator games~\cite{charness2002} gives an advantageous-inequality weight of 0.34, close to the mean of 0.33 across 297 published estimates~\cite{nunnariPozzi2022}, and the asymmetry is significant (likelihood-ratio $\chi^2=24.8$, $p<0.001$).
Social choice therefore shows one reflection axis, self against other, where risky choice was conjectured to have two (Section~\ref{sec:rejected}).

These findings come from work that did not set out to test GDT.
They are consistent with it and do not test it.

%==============================================================================
\section{Four Structural Conjectures Failed}
\label{sec:rejected}
%==============================================================================

Four structural conjectures failed held-out or derivation tests (Table~\ref{tab:rejected}).

\begin{table}[!t]
\renewcommand{\arraystretch}{1.2}
\caption{Conjectures Rejected by Held-Out or Derivation Tests}
\label{tab:rejected}
\centering
\begin{tabular}{@{}p{0.30\columnwidth} p{0.60\columnwidth}@{}}
\toprule
\textbf{Conjecture} & \textbf{How it failed} \\
\midrule
One aversion angle for all domains & Adversarially selected Kahneman--Tversky problem sets shift the angle. It holds only within comparable samples. \\
A conserved quantity of Shannon--Hartley form in the polar encoding & The normalized risk weight changes sign with stakes, which is adaptation and not conservation. \\
Dihedral ($D_4$) symmetry of the fourfold pattern & The apparent chirality comes from the coordinate's endpoints. Gambles at the endpoints are exactly those with no loss branch or no gain branch, verified on 95{,}748 rows, and there the coordinate product reduces to the probability. Fit separately, the pooled chirality of $+0.25\pm0.04$ becomes $-0.03\pm0.07$ inside mixed gambles, on six independent lines. Clustered by problem, the CPC18 estimate is $+0.75$ with a 95\% interval of $[-0.90, +1.38]$. \\
Klein-four symmetry as a general law of risky choice & A registered held-out test on representative gambles scored 0 of 3. The symmetry belongs to the curated problem set. \\
\bottomrule
\end{tabular}
\end{table}

%==============================================================================
\section{Discussion}
\label{sec:discussion}
%==============================================================================

\subsection{What the results establish}
Part I asked whether its metric is a flexible fit or a structure that predicts.
The evidence here favors prediction on a limited set of tests.
On unseen risky-choice problems the model beats expected value and random utility and trails CPT.
Its tradeoff shape transfers between lotteries and games and between two giving tasks.
A low-rank metric frozen before analysis holds on three held-out legs.
The monetary coordinate Part I proposed passes a registered test in three countries.
The model built on Part I's other constants fails that test against a logit, and the temperature rule Part I used fits worse than a fixed temperature on the panel data.

\subsection{What the results do not establish}
The model does not beat CPT on lotteries.
The cross-domain transfer of the full metric is asymmetric, and the symmetric result uses an angle chosen after the fact.
The matched-problem evidence is from language models.
Its five passing contrasts each have an established explanation, and its lottery contrasts mostly fail.
Magnitudes are not calibrated.
The low-rank result holds for gains.
The July predictions were frozen in a local repository, and their order relative to the data rests on commit timestamps.

\subsection{Next tests}
The human matched-problem study under the fixed temperature is designed and frozen (\texttt{prereg-v2}), with sample sizes set from the panel's effect sizes of $d$ between 0.5 and 0.9, and it has not been run.
A contrast that separates GDT from ambiguity aversion, belief change and social preference would need two coordinates manipulated together, so that one metric predicts their joint effect where separate accounts do not.
A registered test of whether one person's risk tradeoff predicts that person's social tradeoff was scored on the Global Preferences Survey~\cite{falk2018} and failed.
Risk correlates with altruism at only $r=0.086$ against a frozen bound of 0.10, and the six preference measures are not low-rank.
That survey reports preference scores rather than choice-fit angles, and a same-subject test on fitted angles is frozen and has no data.
The nine coordinates used here are specified by hand.
A companion study replaces them with directions learned from a sentence encoder and tests those directions on data held out from their discovery (Part III, in preparation).

\subsection{Relevance to computational social systems}
A behavioral component for social simulation is easier to use when its metric has few parameters, transfers from one setting to another, and can be checked with matched problems before an intervention assumes a population ignores a non-monetary attribute.
The panel software is public and runs registered individual-level designs on model populations at scale, which can size a human study before it is run.

%==============================================================================
\section{Conclusion}
\label{sec:conclusion}
%==============================================================================

Part I fit a shared-metric model of economic choice to aggregate targets.
Tested on data it was not fit to, the model beats expected value and random utility on unseen risky-choice problems and trails cumulative prospect theory by a small margin.
A tradeoff shape fit on lotteries predicts play in strategic games, a low-rank metric holds on three held-out data sets, and the income-scaled monetary coordinate Part I proposed passes a registered test.
Matched choice problems move the choices of six language models in the predicted direction on five of seven contrasts, each of which an established account also predicts, and a fixed temperature fits the dose response that Part I's rule did not.
Part I's other constants fail against a logit on high-stakes data, two of seven contrasts fail, and magnitudes are not calibrated.
The human study these results size is designed and has not been run.

%==============================================================================
\section*{Data and Code Availability}
%==============================================================================
\sloppy
Code, frozen predictions and results are in the public repository \url{https://github.com/ahb-sjsu/geometric-economics} under \texttt{experiments/}.
The income-units registration is \texttt{experiments/papers/boundary/} with signed tag \texttt{prereg-boundary-v1} and OSF registration \texttt{osf.io/ead4x}. % confirm OSF public visibility before submission
The panel design and analysis are in \texttt{experiments/projection-gap/} (\texttt{contrasts.py}, \texttt{analyze.py}), and the frozen prediction files of \texttt{prereg-v1} and \texttt{prereg-v2} are \texttt{prereg.lock.json} and \texttt{prereg-v2.lock.json}, public since 2026-07-14.
The \texttt{prereg-sigma-v1} claims are \path{experiments/datasets/PREREG_lowrank_sigma.md}.
The panel harness is the public package \url{https://github.com/ahb-sjsu/llm-panel} (release v0.1.0).
The six model aliases are resolved by the National Research Platform endpoint, and no record maps them to upstream model versions.
CPC18~\cite{cpc18data}, the game data of~\cite{wright2017}, the Andersen et al.\ replication package (openICPSR 112485), and the other data sets are public at the cited sources, and the repository's scripts regenerate every derived file from them.

\section*{Acknowledgment}
Generative AI tools (Claude, Anthropic) were used to assist with analysis code, manuscript drafting, editing, and checking reported numbers against the repository.
The research questions, model, experimental designs, and conclusions are the author's own, and the author is responsible for every reported number.

%==============================================================================
% References
%==============================================================================
\begin{thebibliography}{40}

\bibitem{bond2026partI}
A.~H.~Bond, ``Geometric prediction of economic behavior: Cross-domain validation across game theory and prospect theory,''
\emph{IEEE Trans. Comput. Social Syst.}, 2026, doi: 10.1109/TCSS.2026.3734068. % verify vol/pages on final

\bibitem{bond2026theory}
A.~H.~Bond, ``Testing decision theories by matched designs: Sufficient statistics, sign identification, and scale,'' manuscript in preparation, 2026.

\bibitem{andersen2011}
S.~Andersen, S.~Ertac, U.~Gneezy, M.~Hoffman, and J.~A.~List, ``Stakes matter in ultimatum games,''
\emph{Amer. Econ. Rev.}, vol.~101, no.~7, pp.~3427--3439, Dec. 2011. % verify vol/pages on final

\bibitem{slonim1998}
R.~Slonim and A.~E.~Roth, ``Learning in high stakes ultimatum games: An experiment in the Slovak Republic,''
\emph{Econometrica}, vol.~66, no.~3, pp.~569--596, May 1998. % verify vol/pages on final

\bibitem{cameron1999}
L.~A.~Cameron, ``Raising the stakes in the ultimatum game: Experimental evidence from Indonesia,''
\emph{Econ. Inquiry}, vol.~37, no.~1, pp.~47--59, Jan. 1999. % verify vol/pages on final

\bibitem{erev2017}
I.~Erev, E.~Ert, O.~Plonsky, D.~Roth, and Y.~Roth, ``From anomalies to forecasts: Toward a descriptive model of decisions under risk, under ambiguity, and from experience,''
\emph{Psychol. Rev.}, vol.~124, no.~4, pp.~369--409, 2017.

\bibitem{cpc18data}
O.~Plonsky, I.~Erev, and E.~Ert, ``All raw data for CPC18: Choice prediction competition 2018,''
Zenodo, 2019, doi: 10.5281/zenodo.2571510.

\bibitem{plonsky2019}
O.~Plonsky, R.~Apel, E.~Ert, M.~Tennenholtz, D.~Bourgin, J.~C.~Peterson, D.~Reichman,
T.~L.~Griffiths, S.~J.~Russell, E.~C.~Carter, J.~F.~Cavanagh, and I.~Erev, ``Predicting human
decisions with behavioural theories and machine learning,'' \emph{Nature Hum. Behav.},
vol.~9, pp.~2271--2284, 2025, doi: 10.1038/s41562-025-02267-6.

\bibitem{wright2017}
J.~R.~Wright and K.~Leyton-Brown, ``Predicting human behavior in unrepeated, simultaneous-move games,''
\emph{Games Econ. Behav.}, vol.~106, pp.~16--37, Nov. 2017.

\bibitem{fehr1999}
E.~Fehr and K.~M.~Schmidt, ``A theory of fairness, competition, and cooperation,''
\emph{Quart. J. Econ.}, vol.~114, no.~3, pp.~817--868, Aug. 1999.

\bibitem{fraser2020}
B.~J.~Fraser and D.~Nettle, ``Hunger affects social decisions in a multi-round public goods game but not an ultimatum game,''
\emph{Adaptive Human Behav. Physiol.}, vol.~6, pp.~334--355, 2020.

\bibitem{ruggeri2020}
K.~Ruggeri \emph{et~al.}, ``Replicating patterns of prospect theory for decision under risk,''
\emph{Nature Human Behav.}, vol.~4, no.~6, pp.~622--633, Jun. 2020.

\bibitem{fisman2007}
R.~Fisman, S.~Kariv, and D.~Markovits, ``Individual preferences for giving,''
\emph{Amer. Econ. Rev.}, vol.~97, no.~5, pp.~1858--1876, Dec. 2007.

\bibitem{charness2002}
G.~Charness and M.~Rabin, ``Understanding social preferences with simple tests,''
\emph{Quart. J. Econ.}, vol.~117, no.~3, pp.~817--869, Aug. 2002.

\bibitem{ellsberg1961}
D.~Ellsberg, ``Risk, ambiguity, and the Savage axioms,''
\emph{Quart. J. Econ.}, vol.~75, no.~4, pp.~643--669, Nov. 1961.

\bibitem{andreoni1990}
J.~Andreoni, ``Impure altruism and donations to public goods: A theory of warm-glow giving,''
\emph{Econ. J.}, vol.~100, no.~401, pp.~464--477, Jun. 1990.

\bibitem{akerlof2000}
G.~A.~Akerlof and R.~E.~Kranton, ``Economics and identity,''
\emph{Quart. J. Econ.}, vol.~115, no.~3, pp.~715--753, Aug. 2000.

\bibitem{falk2018}
A.~Falk, A.~Becker, T.~Dohmen, B.~Enke, D.~Huffman, and U.~Sunde, ``Global evidence on economic preferences,''
\emph{Quart. J. Econ.}, vol.~133, no.~4, pp.~1645--1692, Nov. 2018. % verify vol/pages on final

\bibitem{sims2003}
C.~A.~Sims, ``Implications of rational inattention,''
\emph{J. Monetary Econ.}, vol.~50, no.~3, pp.~665--690, Apr. 2003.

\bibitem{matejka2015}
F.~Mat\v{e}jka and A.~McKay, ``Rational inattention to discrete choices: A new foundation for the multinomial logit model,''
\emph{Amer. Econ. Rev.}, vol.~105, no.~1, pp.~272--298, Jan. 2015.

\bibitem{bruhin2019}
A.~Bruhin, E.~Fehr, and D.~Schunk, ``The many faces of human sociality: Uncovering the distribution and stability of social preferences,''
\emph{J. Eur. Econ. Assoc.}, vol.~17, no.~4, pp.~1025--1069, Aug. 2019.

\bibitem{fehrCharness2025}
E.~Fehr and G.~Charness, ``Social preferences: Fundamental characteristics and economic consequences,''
\emph{J. Econ. Literature}, vol.~63, no.~2, pp.~440--514, Jun. 2025.

\bibitem{nunnariPozzi2022}
S.~Nunnari and M.~Pozzi, ``Meta-analysis of inequality aversion estimates,''
CESifo Working Paper No.~9851, Munich, Germany, 2022.

\end{thebibliography}

\balance

\end{document}
